
Practitioners implicitly know to avoid horizontal flip augmentation when training models on MNIST digits---Kaggle winners exclude it from winning solutions, PyTorch official tutorials omit it---yet no formal validation explains why, or quantifies the harm. This gap between implicit knowledge and rigorous understanding reflects a broader problem in data augmentation practice: standard transformations are assumed ''safe'' without domain-specific validation, potentially introducing label noise that degrades accuracy on affected classes.

The consequences can be serious in safety-critical domains. A chest X-ray model trained with flip augmentation may silently degrade on lateralized pathologies (left versus right lung conditions), where anatomical orientation is semantically critical. Traffic sign recognition systems trained with flip on directional arrows may misclassify critical navigational signs. In these domains, undetected semantic violations could lead to silent failures with real-world consequences masked by aggregate accuracy metrics.

Data augmentation has become a cornerstone technique for improving deep learning generalization by artificially expanding training sets with label-preserving transformations. Standard augmentations---horizontal flip, rotation, crop, color jitter---are catalogued in comprehensive surveys and deployed widely across computer vision tasks. However, augmentations that algorithmically preserve labels may violate domain-specific semantic constraints. Consider horizontal flip applied to MNIST: flipping an asymmetric digit like '2' produces a visually non-canonical image that retains its original label. The model sees both canonical and mirror-reversed versions of '2' during training, yet test sets contain only canonical digits. This creates a mismatch between training augmentation invariances and the domain's actual symmetry structure. When such misalignment occurs, augmentation transforms from a regularization technique into a source of systematic label noise.

Why has this problem gone unaddressed? First, practitioners often rely on aggregate metrics that obscure class-specific degradation. Second, augmentation surveys focus on cataloguing techniques and their general benefits rather than analyzing semantic validity constraints. Third, the label noise literature concentrates on annotation errors but does not consider augmentation-induced noise as a distinct category. Finally, semantic validity is treated as implicit practitioner knowledge encoded in folklore rather than an explicit, testable design criterion.

We reframe this problem through a mechanistic lens: augmentation effectiveness depends on alignment between augmentation-induced invariances and domain-specific symmetries. Horizontal flip on MNIST creates label noise because asymmetric digits lack horizontal reflection symmetry---flipped images are semantically invalid for their retained labels. This insight is testable through controlled experiments: if semantic invalidity causes degradation, we should observe (1) differential effects on asymmetric versus symmetric digits, (2) dose-response relationships where degradation magnitude increases with flip probability, and (3) absence of effects under semantically valid augmentations like rotation.

We formalize and validate this semantic validity hypothesis through rigorous experimentation on MNIST digit classification. Our core finding is striking: horizontal flip introduces label noise for asymmetric digits {2, 3, 5, 6, 7, 9}, causing statistically significant test accuracy degradation of 0.72-1.00 percentage points at moderate flip probability (p=0.5), with dose-dependent degradation ranging from 0.37 pp (p=0.3) to 4.10 pp (p=0.9) while symmetric digits {0, 1, 8} remain largely stable at moderate flip rates. The degradation exhibits perfect or near-perfect dose-response relationships (Spearman $\rho$ = -1.0 in h-m1, $\rho$ = -0.969 in h-m), revealing a deterministic mechanism. Critically, rotation $\pm 15$° augmentation shows no differential effect across three independent validations, isolating semantic invalidity as the causal factor.

Our contributions are:

\begin{enumerate}
\item \textbf{First rigorous semantic validity test on MNIST}: We provide the first controlled experimental validation of semantic validity for a standard data augmentation technique, testing horizontal flip on MNIST with four independent sub-hypotheses spanning existence, mechanism, and condition validation.
\end{enumerate}

\begin{enumerate}
\item \textbf{Quantification of augmentation-induced label noise}: We demonstrate that flip probability directly controls label noise rate, enabling systematic measurement of effect sizes ranging from -0.37\% (flip probability 0.3) to -4.10\% (flip probability 0.9) on asymmetric digit accuracy.
\end{enumerate}

\begin{enumerate}
\item \textbf{Perfect dose-response evidence}: We establish a dose-response framework for augmentation validation, achieving Spearman $\rho$ = -1.0 in one experiment and $\rho$ = -0.969 in another, indicating a deterministic relationship between flip rate and degradation magnitude.
\end{enumerate}

\begin{enumerate}
\item \textbf{Semantic validity framework demonstrated on MNIST}: We demonstrate semantic validity testing on MNIST digits and propose a generalizable framework: an augmentation is valid if and only if augmented images remain semantically in-distribution for their labeled class. We operationalize this through differential effect testing (asymmetric vs. symmetric digits), dose-response analysis, and positive control validation (rotation vs. flip).
\end{enumerate}

This work bridges data augmentation surveys that catalogue techniques without semantic analysis and label noise research that overlooks augmentation as a systematic noise source. By converting practitioner folklore into validated science for MNIST, we establish actionable guidelines grounded in quantitative evidence. More broadly, we demonstrate that augmentation design requires explicit validation against domain-specific semantic constraints.

